\chapter{Scabies}
\index{Scabies}
\label{sec:Scabies}

\section{Overview}


\begin{itemize}[noitemsep]
  \item Scabies is a parasitic infestation of the skin caused by the mite \textit{Sarcoptes scabiei var.~hominis}. Transmission occurs through prolonged skin-to-skin contact
  \item Fomite transmission is less common but possible.
\end{itemize}
\section{Causes}
This condition happens because of the female mite burrowing into the stratum corneum, laying eggs and triggering a delayed type IV hypersensitivity reaction.     
\section{Risk Factors}

\begin{itemize}[noitemsep]
  \item Genetic predisposition and family history
  \item Advancing age
  \item Lifestyle factors (diet, exercise, smoking, alcohol)
  \item Environmental exposures
  \item Underlying medical conditions
  \item Medication use
\end{itemize}
\section{Signs and Symptoms}

\subsection{Common symptoms}
\begin{itemize}[noitemsep]
  \item You may experience intensely pruritic (especially nocturnal), papular, and vesicular eruptions in characteristic locations: interdigital web spaces, wrists, elbows, axillae, periumbilical area, waistband, buttocks, and genitalia in males
  \item Burrows (thin, linear, grayish-white tracks) are pathognomonic. Crusted (Norwegian) scabies is a highly contagious variant in immunocompromised patients with widespread, thick, hyperkeratotic crusts.
  \item Pain or discomfort in the affected area
  \end{itemize}

\subsection{Emergency warning signs}
\begin{itemize}[noitemsep]
  \item Severe or worsening symptoms of scabies require immediate medical evaluation
\end{itemize}
\section{Progression and Stages}
The condition may progress through different stages of severity over time. Early stages often involve mild or subtle symptoms that may be overlooked. Moderate stages involve more noticeable symptoms affecting daily function. Advanced stages may involve significant impairment and complications.
\section{Complications}
\begin{itemize}[noitemsep]
  \item Disease progression leading to worsening symptoms
  \item Secondary infections or injuries
  \item Side effects of treatment
  \item Reduced quality of life
  \item Emotional and psychological impact
\end{itemize}
\section{Diagnosis}
Doctors diagnose this based on clinical examination

\begin{itemize}[noitemsep]
  \item Dermoscopy may reveal the "delta wing jet" appearance
  \item Skin scraping confirms.
  \item Comprehensive medical history and physical examination
  \item Laboratory tests to assess organ function
  \item Imaging studies to visualize affected structures
  \item Specialized testing depending on the affected system
  \item Consultation with relevant specialists
\end{itemize}
\section{Prevention}
\begin{itemize}[noitemsep]
  \item Maintain a healthy lifestyle with balanced diet and exercise
  \item Avoid known risk factors
  \item Attend regular medical check-ups for early detection
  \item Follow recommended screening guidelines
\end{itemize}
\section{Treatment Options}
\begin{itemize}[noitemsep]
  \item Treatment typically involves with topical permethrin 5\% cream (applied to the entire body from neck down, left on for 8--14 hours, repeated in 1 week) or oral ivermectin (200 $\mu$g/kg, repeated in 1--2 weeks)
  \item Close contacts should be treated simultaneously.
  \item Medications to manage symptoms and address underlying mechanisms
  \item Lifestyle modifications and self-care measures
  \item Physical therapy and rehabilitation
  \item Surgical intervention when indicated
  \item Regular monitoring and follow-up care
\end{itemize}
\section{Living with It}
\begin{itemize}[noitemsep]
  \item Follow your treatment plan as prescribed
  \item Maintain regular follow-up with your healthcare provider
  \item Make necessary lifestyle adjustments
  \item Seek support from family, friends, or support groups
  \item Stay informed about your condition
\end{itemize}
\section{Outlook}
\begin{itemize}[noitemsep]
  \item Most people recover well with treatment
  \item Post-scabetic itching may persist for weeks.
\end{itemize}
